
\subsubsection{2.1 Output-Level Uncertainty Quantification}

Token-based methods measure uncertainty from output distributions. Fadeeva et al. (2024) survey token-level uncertainty metrics including entropy, perplexity, and sequence probability. These methods achieve 0.62--0.65 AUROC for correctness prediction on factual QA. The fundamental limitation is that output softmax distributions conflate linguistic confidence (fluency) with factual confidence (correctness).

Semantic entropy \citep{kuhn2023semantic} addresses this by sampling multiple responses and clustering by semantic equivalence, achieving approximately 0.80 AUROC. However, this requires 5--20 forward passes per question.

Kossen et al. (2024) train probes to predict semantic entropy from hidden states, achieving correlation with ground-truth semantic entropy. However, predicting uncertainty differs from predicting correctness: a confident model can still be wrong.

\subsubsection{2.2 Hidden State Probing}

Linear probes have successfully extracted semantic concepts from transformer hidden states, including sentiment, syntax, and factual knowledge. Burns et al. (2023) use contrastive probing to detect model beliefs. Our work extends this paradigm: rather than probing for specific facts, we probe for correctness of model-generated answers.

Semantic Entropy Probes \citep{kossen2024semantic} train linear probes on hidden states to predict semantic entropy. We depart from this approach in two ways: (1) our target is ground-truth correctness, not uncertainty; (2) we systematically characterize layer depth effects, identifying optimal extraction at 50\% depth.

\subsubsection{2.3 Transformer Interpretability}

The logit lens (nostalgebraist, 2020) demonstrates that projecting intermediate hidden states to vocabulary space reveals semantic content developing across layers. This supports the hypothesis that middle layers encode semantic knowledge before output formatting.

Aiersilan et al. (2026) study hallucination detection across layers, finding optimal signals at 40--56\% depth for Llama and Mistral models. Our inverted-U pattern (peak at 50\% depth) aligns with these findings.

